% =============================================================
%  § 3.  Криптосистемы RSA и Диффи--Хеллмана
% =============================================================

\section{Криптосистемы RSA и Диффи--Хеллмана}
\label{sec:rsa-dh}

В этом параграфе мы построим две настоящие, работающие сегодня схемы асимметричной криптографии. Обе появились почти в одно время --- между 1976 и 1978 годами --- и обе опираются на результаты предыдущего параграфа.

\subsection{Криптосистема RSA}

Идея была опубликована в 1978 году тремя сотрудниками Массачусетского технологического института --- \emph{Р}ональдом Райвестом, \emph{Ш}амиром (Ади Шамир) и \emph{А}длеманом (Леонард Адлеман). Их фамилии и дали название системе --- RSA.

\paragraph*{Генерация ключей.} Каждый пользователь, желающий получать зашифрованные сообщения, выполняет одноразовую процедуру.

\begin{algorithmbox}[title={Алгоритм: генерация ключей RSA}]
\begin{enumerate}
  \item Выбрать два больших случайных \emph{различных} простых числа $p$ и $q$ (на практике --- по 1024 бита каждое).
  \item Вычислить $n = p\cdot q$ и $\Eul(n) = (p-1)(q-1)$.
  \item Выбрать целое $e$, взаимно простое с $\Eul(n)$ (часто берут $e = 65537$).
  \item С помощью расширенного алгоритма Евклида найти $d$ такое, что
  \[
  e\cdot d \equiv 1 \modn{\Eul(n)}.
  \]
  \item \term{Открытый ключ:} пара $(n, e)$ --- публикуется.\\
        \term{Закрытый ключ:} число $d$ --- хранится в секрете.\\
        \term{Лишние данные:} числа $p, q, \Eul(n)$ безопасно уничтожаются.
\end{enumerate}
\end{algorithmbox}

\paragraph*{Шифрование.} Алиса хочет послать сообщение $m$ Бобу. Сообщение представлено числом из отрезка $[0, n-1]$ (длинное сообщение разбивается на блоки). Алиса берёт открытый ключ Боба $(n, e)$ и вычисляет:
\[
c \;=\; m^e \bmod n.
\]
Число $c$ называется \term{шифр-текстом} и передаётся открыто.

\paragraph*{Расшифрование.} Боб, получив $c$, использует свой закрытый ключ $d$ и вычисляет:
\[
m' \;=\; c^d \bmod n.
\]
Утверждается: $m' = m$ (рис.~\ref{fig:rsa-scheme}).

\begin{figure}[H]
\centering
\begin{tikzpicture}[
  font=\small,
  every node/.style={align=center},
  box/.style={rectangle, draw=prosvBlue, fill=prosvLight,
              minimum width=3.0cm, minimum height=1.1cm,
              rounded corners=2pt, font=\small}
]
  \node[circle, draw=prosvBlue, fill=white, minimum size=0.8cm] (a) at (-5.5,0) {A};
  \node[below, font=\scriptsize] at (-5.5,-0.55) {Алиса};
  \node[circle, draw=prosvBlue, fill=white, minimum size=0.8cm] (b) at (5.5,0) {B};
  \node[below, font=\scriptsize] at (5.5,-0.55) {Боб};
  
  \node[box] (enc) at (-2.5,0) {$c = m^e \bmod n$};
  \node[box] (dec) at (2.5,0) {$m = c^d \bmod n$};
  
  \draw[->,>=stealth, prosvBlue, thick] (a) -- node[above]{$m$} (enc);
  \draw[->,>=stealth, prosvBlue, thick] (enc) -- node[above]{$c$} (dec);
  \draw[->,>=stealth, prosvBlue, thick] (dec) -- node[above]{$m$} (b);
  
  % Ключи
  \node[font=\scriptsize, text=prosvGreen, align=center] at (-2.5, 1.4) {откр. ключ Боба\\$(n, e)$};
  \node[font=\scriptsize, text=prosvRed, align=center] at (2.5, 1.4) {закр. ключ Боба\\$d$};
  
  \draw[->, dashed, prosvGreen] (-2.5, 1.0) -- (-2.5, 0.6);
  \draw[->, dashed, prosvRed] (2.5, 1.0) -- (2.5, 0.6);
  
  \draw[->,>=stealth, prosvGreen, dashed] (b) to[bend right=50] (-2.5, 1.6);
  \node[font=\scriptsize, text=prosvGreen] at (1, 2.4) {опубликован};
\end{tikzpicture}
\caption{Схема работы RSA: шифрование --- возведение в степень $e$ по модулю $n$, расшифрование --- возведение в степень $d$. Обе операции эффективны благодаря быстрому возведению в степень.}
\label{fig:rsa-scheme}
\end{figure}

% -----------------------------------------------------------------
\subsection{Почему RSA работает: корректность}

\begin{theorem}{корректность RSA}{thm:p3-1}
В обозначениях процедуры RSA, для любого $m \in [0, n-1]$ выполняется
\[
(m^e)^d \equiv m \modn{n}.
\]
\end{theorem}

\begin{proof}
По построению $e\cdot d \equiv 1 \modn{\Eul(n)}$, значит, найдётся целое $k$, такое что $e\cdot d = 1 + k\cdot \Eul(n)$. Подставим:
\[
m^{ed} = m^{1 + k\Eul(n)} = m\cdot \bigl(m^{\Eul(n)}\bigr)^k.
\]

\textbf{Случай 1.} $\NOD(m, n) = 1$. Тогда по теореме Эйлера $m^{\Eul(n)} \equiv 1 \modn n$, и
\[
m^{ed} \equiv m\cdot 1^k = m \modn n.
\]

\textbf{Случай 2.} $\NOD(m, n) > 1$. Поскольку $n = pq$ и $p, q$ простые, это значит, что $m$ делится либо на $p$, либо на $q$ (но не на оба --- иначе $m \ge pq = n$, чего быть не может). Рассмотрим, например, $p \mid m$. Тогда $m \equiv 0 \modn p$, и обе части сравнения $m^{ed} \equiv m \modn p$ обращаются в ноль --- утверждение выполнено.

С другой стороны, $\NOD(m, q) = 1$, поэтому по малой теореме Ферма $m^{q-1} \equiv 1 \modn q$. А значит, $m^{(p-1)(q-1)} = m^{\Eul(n)} \equiv 1 \modn q$, и
\[
m^{ed} = m\cdot \bigl(m^{\Eul(n)}\bigr)^k \equiv m\cdot 1^k = m \modn q.
\]
Получили $m^{ed} \equiv m$ одновременно по модулям $p$ и $q$. По китайской теореме об остатках (или просто потому, что разность $m^{ed} - m$ делится и на $p$, и на $q$, а значит, на их произведение $n$) сравнение $m^{ed} \equiv m \modn n$ выполнено. \qedhere
\end{proof}

\paragraph*{Учебный пример (маленькие числа).} Полноценный RSA использует тысячебитные числа, но для понимания работает и крошечный пример.

\begin{example}{генерация и применение ключа RSA}{ex:p3-1}
Пусть $p = 11$, $q = 23$. Тогда:
\begin{itemize}
  \item $n = 11\cdot 23 = 253$, $\Eul(n) = 10\cdot 22 = 220$.
  \item Выберем $e = 7$. Проверим взаимную простоту: $\NOD(7, 220) = 1$.
  \item Найдём $d = 7^{-1} \bmod 220$. Расширенным алгоритмом Евклида:
  \[
  220 = 31\cdot 7 + 3,\quad 7 = 2\cdot 3 + 1.
  \]
  Раскручиваем: $1 = 7 - 2\cdot 3 = 7 - 2(220 - 31\cdot 7) = 63\cdot 7 - 2\cdot 220$. Значит, $d = 63$ (проверка: $7\cdot 63 = 441 = 2\cdot 220 + 1$).
  \item Открытый ключ: $(253, 7)$. Закрытый ключ: $d = 63$.
\end{itemize}

\emph{Шифрование.} Пусть $m = 50$. Считаем $c = 50^7 \bmod 253$ быстрым возведением:
\[
50^2 = 2500 \equiv 2500 - 9\cdot 253 = 223,\qquad 50^4 \equiv 223^2 = 49\,729.
\]
Делим $49\,729 / 253 \approx 196{,}5$, $196\cdot 253 = 49\,588$, остаток $141$. Значит, $50^4 \equiv 141$. Далее $50^7 = 50^4\cdot 50^2\cdot 50 \equiv 141\cdot 223\cdot 50 \bmod 253$. Считаем: $141\cdot 223 = 31\,443 = 124\cdot 253 + 71$, значит $\equiv 71$. Затем $71\cdot 50 = 3550 = 14\cdot 253 + 8$, значит $\equiv 8$. Итак, $c = 8$.

\emph{Расшифрование.} Боб считает $c^d = 8^{63} \bmod 253$. Вместо длинных вычислений «вручную» можно использовать малую теорему Ферма по модулям 11 и 23 раздельно, а затем китайскую теорему об остатках --- именно так и поступают на практике. Ответ: $50$.
\end{example}

\begin{interestbox}
Удивительный исторический парадокс: ровно та же схема была независимо открыта в 1973 году британским математиком Клиффордом Коксом, работавшим в Британском управлении правительственной связи (GCHQ). Но работа была засекречена и рассекречена только в 1997 году. К этому моменту Райвест, Шамир и Адлеман уже стали всемирно известными, основали компанию RSA Security, а Шамир получил Премию Тьюринга. Типичный пример того, как секретность тормозит развитие науки.
\end{interestbox}

% -----------------------------------------------------------------
\subsection{Почему RSA безопасен}

\begin{importantbox}
Стойкость RSA опирается на предположение: \emph{задача факторизации $n$ практически нерешаема при больших $n$}.
\end{importantbox}

Действительно, если злоумышленник умеет факторизовать $n$, он восстанавливает $p, q$, затем $\Eul(n) = (p-1)(q-1)$ и наконец $d = e^{-1} \bmod \Eul(n)$ --- получит закрытый ключ. Обратное (что взлом RSA эквивалентен факторизации) также считается верным, но строго не доказано: это одна из открытых проблем в криптографии.

Заметим, что \emph{знать $n$ и $e$} (открытый ключ) --- ещё далеко не значит знать $\Eul(n)$. Знание $\Eul(n)$ и знание разложения $n = pq$ --- одна и та же информация: ведь $p + q = n - \Eul(n) + 1$, а $p - q = \sqrt{(p+q)^2 - 4n}$. То есть «нахождение $\Eul(n)$» и «факторизация $n$» --- задачи одинаковой сложности.

\paragraph*{Почему нужны именно \emph{большие} простые числа?} Современный рекомендуемый размер $n$ --- 2048 или даже 4096 бит. Чтобы взломать 2048-битный RSA, нужно факторизовать число длиной около 617 десятичных цифр; лучший известный алгоритм потратит на это больше времени, чем существует Вселенная.

\begin{example}{границы практической факторизации}{ex:p3-2}
В 2020 году группе исследователей удалось факторизовать число RSA-250 (250 десятичных цифр $\approx$ 829 бит). Заняло это около 2700 ядро-лет вычислений --- то есть один компьютер с 2700 ядрами работал бы целый год. Это подтверждает, что 1024-битный RSA скоро будет «на грани», а 2048-битный --- надолго безопасен.
\end{example}

\paragraph*{Проверим выполнение требований из § \ref{sec:history}.}
\begin{enumerate}
  \item \emph{Корректность} --- доказана теоремой выше.
  \item \emph{Стойкость} --- опирается на сложность факторизации.
  \item \emph{Эффективность} --- шифрование/расшифрование сводятся к возведению в степень по модулю, что мы умеем делать полиномиально (см. §~\ref{subsec:fastpow}).
  \item \emph{Защита ключа} --- знание $e, n$ не даёт быстрого способа найти $d$ без факторизации.
\end{enumerate}

% -----------------------------------------------------------------
\subsection{Протокол Диффи--Хеллмана}

Уитфилд Диффи и Мартин Хеллман решали несколько другую (хотя и родственную) задачу: \emph{договориться} о секретном ключе по открытому каналу, --- не передавая по нему ничего секретного. Их работа 1976 года так и называется: «Новые направления в криптографии».

\paragraph*{Идея на пальцах: краски.} Прежде чем вводить формулы, разберём наглядную аналогию (рис.~\ref{fig:p3-2}).

\begin{figure}[H]
\centering
\begin{tikzpicture}[
  font=\small, every node/.style={align=center},
  circ/.style={circle, draw, minimum size=0.7cm}
]
  % Алиса и Боб + общая краска
  \node[circ, fill=yellow!50!brown!20] (g) at (0,3) {};
  \node[below, font=\scriptsize] at (0,2.6) {общая «жёлтая»};
  
  \node[circ, fill=red!60!brown!50] (sa) at (-5,3) {};
  \node[below, font=\scriptsize, align=center] at (-5,2.6) {Алисин «красный»\\(секрет)};
  
  \node[circ, fill=blue!60!cyan!50] (sb) at (5,3) {};
  \node[below, font=\scriptsize, align=center] at (5,2.6) {Бобов «синий»\\(секрет)};
  
  % Смеси
  \node[circ, fill=orange!70!red!30] (ya) at (-2.5,1.0) {};
  \node[below, font=\scriptsize, align=center] at (-2.5,0.6) {жёлт. $+$ красн.\\(передаётся\\открыто)};
  
  \node[circ, fill=green!60!blue!30] (yb) at (2.5,1.0) {};
  \node[below, font=\scriptsize, align=center] at (2.5,0.6) {жёлт. $+$ син.\\(передаётся\\открыто)};
  
  \draw[->, prosvBlue] (g) -- (ya);
  \draw[->, prosvBlue] (sa) -- (ya);
  \draw[->, prosvBlue] (g) -- (yb);
  \draw[->, prosvBlue] (sb) -- (yb);
  
  % Итоговая общая смесь
  \node[circ, fill=brown!50] (final) at (0,-2.0) {};
  \node[below, font=\scriptsize, align=center] at (0,-2.4) {итоговая «коричневая»\\ --- общий секрет};
  
  \draw[->, prosvGreen, thick] (yb) -- node[above right, font=\scriptsize, text=prosvGreen]{+ Алисин «красный»} (final);
  \draw[->, prosvGreen, thick] (ya) -- node[above left, font=\scriptsize, text=prosvGreen]{+ Бобов «синий»} (final);
\end{tikzpicture}
\caption{Аналогия Диффи--Хеллмана со смешиванием красок. Смешать --- легко, разделить смесь обратно на исходные цвета --- невозможно. Алиса и Боб приходят к одному цвету разными путями.}
\label{fig:p3-2}
\end{figure}

\noindent
Алиса и Боб условились о «жёлтой» базовой краске (известна всем). Каждый выбрал свою секретную краску (Алиса --- красный, Боб --- синий) и публично передал смесь «своя $+$ жёлтая». Подслушивающий видит обе смеси, но не может разложить их обратно. А Алиса и Боб, добавив к чужой смеси свою секретную краску, приходят к одинаковому итогу: жёлт. $+$ красн. $+$ син.

\paragraph*{Математическая реализация.} Роль «смешивания» играет \term{дискретное возведение в степень} по модулю простого числа (рис.~\ref{fig:dh-scheme}).

\begin{algorithmbox}[title={Алгоритм: протокол Диффи--Хеллмана}]
\textbf{Подготовка (общедоступно):} большое простое число $p$ и целое $g$, $1 < g < p$.
\begin{enumerate}
  \item Алиса выбирает случайный секрет $a \in [1, p-2]$ и посылает Бобу $A = g^a \bmod p$.
  \item Боб выбирает случайный секрет $b \in [1, p-2]$ и посылает Алисе $B = g^b \bmod p$.
  \item Алиса вычисляет общий ключ: $K = B^a \bmod p$.
  \item Боб вычисляет общий ключ: $K = A^b \bmod p$.
\end{enumerate}
\textbf{Результат.} Алиса и Боб получили один и тот же ключ
\[
K = g^{ab} \bmod p,
\]
который никогда не передавался по каналу связи.
\end{algorithmbox}

\begin{figure}[H]
\centering
\begin{tikzpicture}[
  font=\small,
  every node/.style={align=center},
  alice/.style={fill=prosvLight, draw=prosvBlue, rounded corners=2pt, minimum width=4.5cm, minimum height=0.7cm},
  bob/.style={fill=prosvCream, draw=prosvGold!80!brown, rounded corners=2pt, minimum width=4.5cm, minimum height=0.7cm}
]
  % Заголовки
  \node[font=\bfseries\color{prosvBlue}] at (-3.5, 4) {Алиса};
  \node[font=\bfseries\color{prosvRed}] at (3.5, 4) {Боб};
  \node[font=\bfseries\color{prosvGreen}] at (0, 4) {Открытый канал};
  
  % Шаги Алисы
  \node[alice] (A1) at (-3.5, 3) {Выбирает секрет $a$};
  \node[alice] (A2) at (-3.5, 2) {Считает $A = g^a \bmod p$};
  \node[alice] (A3) at (-3.5, 0.5) {Получает $B$};
  \node[alice] (A4) at (-3.5, -0.5) {$K = B^a \bmod p$};
  
  % Шаги Боба
  \node[bob] (B1) at (3.5, 3) {Выбирает секрет $b$};
  \node[bob] (B2) at (3.5, 2) {Считает $B = g^b \bmod p$};
  \node[bob] (B3) at (3.5, 0.5) {Получает $A$};
  \node[bob] (B4) at (3.5, -0.5) {$K = A^b \bmod p$};
  
  % Передача
  \draw[->, thick, prosvBlue] (A2) -- node[above, font=\scriptsize]{$A$} (B3);
  \draw[->, thick, prosvRed]  (B2) -- node[below, font=\scriptsize]{$B$} (A3);
  
  % Итог
  \node[fill=prosvLight!30!white, draw=prosvGreen, very thick, rounded corners=3pt, minimum width=7cm, minimum height=0.8cm] at (0,-1.8) {Общий ключ $K = g^{ab} \bmod p$};
\end{tikzpicture}
\caption{Диаграмма обмена в протоколе Диффи--Хеллмана.}
\label{fig:dh-scheme}
\end{figure}

\paragraph*{Учебный пример.}

\begin{example}{обмен ключами Диффи--Хеллмана}{ex:p3-3}
Пусть $p = 23$, $g = 5$.
\begin{itemize}
  \item Алиса выбирает $a = 6$, отправляет $A = 5^6 \bmod 23$. Считая последовательно по таблице степеней пятёрки:
  \[
  5^1\equiv 5,\;5^2\equiv 2,\;5^3\equiv 10,\;5^4\equiv 4,\;5^5\equiv 20,\;5^6\equiv 8,
  \]
  получаем $A = 8$.
  \item Боб выбирает $b = 15$, отправляет $B = 5^{15} \bmod 23$. Продолжая таблицу:
  \[
  5^7\equiv 17,\;5^8\equiv 16,\;5^9\equiv 11,\;5^{10}\equiv 9,\;5^{11}\equiv 22,
  \]
  \[
  5^{12}\equiv 18,\;5^{13}\equiv 21,\;5^{14}\equiv 13,\;5^{15}\equiv 19,
  \]
  получаем $B = 19$.
  \item Алиса вычисляет общий ключ $K = B^a = 19^6 \bmod 23$. Замечая, что $19 \equiv -4 \modn{23}$, имеем $(-4)^6 = 4^6 = 4096 = 178\cdot 23 + 2$, значит, $K = 2$.
  \item Боб вычисляет $K = A^b = 8^{15} \bmod 23$. Но $8 = 5^6$, значит, $8^{15} = 5^{90}$. По малой теореме Ферма $5^{22} \equiv 1 \modn{23}$, а $90 = 4\cdot 22 + 2$, поэтому $5^{90} \equiv 5^2 = 2 \modn{23}$. Получаем $K = 2$.
\end{itemize}

Согласование удалось: общий ключ $K = 2$. И это притом, что числа $a = 6$ и $b = 15$ никогда не передавались по каналу!
\end{example}

\paragraph*{Безопасность Диффи--Хеллмана.} Чтобы подслушивающий, видя $g, p, A, B$, мог найти общий ключ $K = g^{ab}$, ему нужно сначала восстановить хотя бы один из секретов $a$ или $b$, --- то есть решить уравнение
\[
g^a \equiv A \modn p
\]
относительно $a$ при известных $g, A, p$. Эта задача называется \term{задачей дискретного логарифма} (в группе $\Zn{p}^*$).

\begin{importantbox}
Для всех известных классических алгоритмов задача дискретного логарифма по модулю большого простого --- субэкспоненциальная (примерно той же сложности, что и факторизация). Однако \emph{по модулю эллиптической кривой} --- значительно сложнее, и поэтому современные мессенджеры обычно используют вариант Диффи--Хеллмана не по модулю числа, а на эллиптических кривых (\engterm{ECDH}). Идея та же.
\end{importantbox}

% -----------------------------------------------------------------
\subsection{Атака «человек посередине»}

И RSA, и Диффи--Хеллман в чистом виде уязвимы к одной атаке. Представим, что между Алисой и Бобом «вклинился» злоумышленник --- Мэллори (\engterm{Man-in-the-Middle}) (рис.~\ref{fig:p3-4}).

\begin{figure}[H]
\centering
\begin{tikzpicture}[font=\small, every node/.style={align=center}]
  \node[circle, draw=prosvBlue, fill=white, minimum size=0.8cm] (a) at (-5,0) {A};
  \node[below, font=\scriptsize] at (-5,-0.55) {Алиса};
  
  \node[circle, draw=prosvRed, fill=prosvCream, minimum size=0.8cm] (m) at (0,0) {M};
  \node[below, font=\scriptsize, text=prosvRed] at (0,-0.55) {Мэллори};
  
  \node[circle, draw=prosvBlue, fill=white, minimum size=0.8cm] (b) at (5,0) {B};
  \node[below, font=\scriptsize] at (5,-0.55) {Боб};
  
  \draw[->, prosvBlue, thick] (a) -- node[above, font=\scriptsize]{«Боб, мой $A$»} (m);
  \draw[->, prosvRed, thick]  (m) -- node[above, font=\scriptsize]{«Я Алиса, мой $A'$»} (b);
  \draw[<-, prosvBlue, thick] (a) to[bend right=30] node[below, font=\scriptsize]{«моё $B''$»} (m);
  \draw[<-, prosvRed, thick]  (m) to[bend right=30] node[below, font=\scriptsize]{«моё $B$»} (b);
\end{tikzpicture}
\caption{Атака «человек посередине»: Мэллори создаёт две независимые сессии, в каждой выдавая себя за противоположную сторону.}
\label{fig:p3-4}
\end{figure}

Мэллори перехватывает $A$ от Алисы, заменяет на свой $A'$ и пересылает Бобу. С Бобом он устанавливает один ключ, с Алисой --- другой. Затем расшифровывает каждое сообщение, читает (и при желании меняет), снова шифрует и передаёт дальше. Внешне Алиса и Боб ничего не замечают.

\paragraph*{Защита.} Нужна \emph{аутентификация} --- способ убедиться, что собеседник действительно тот, за кого себя выдаёт. В реальных системах это решается через \term{цифровые сертификаты} и \term{электронную цифровую подпись}. О них --- следующий параграф.

% -----------------------------------------------------------------
\begin{tasksbox}
\begin{enumerate}
  \item В системе RSA дано $p = 7$, $q = 13$, $e = 5$. Найдите $n$, $\Eul(n)$ и $d$.
  \item Используя ключ из задачи 1, зашифруйте сообщение $m = 9$.
  \item В Диффи--Хеллмане положите $p = 17$, $g = 3$, $a = 5$, $b = 7$. Найдите $A$, $B$ и общий ключ $K$.
  \item Почему в RSA нельзя брать $e = 1$? А $e = \Eul(n) - 1$?
  \item Что произойдёт, если в RSA Алиса случайно использует $m \ge n$?
  \item* Объясните, как именно атака «человек посередине» работает в случае RSA (не Диффи--Хеллмана).
  \item* Что вы посоветуете человеку, который хочет «шифровать всё своими силами без сертификатов»? Какие угрозы он не учитывает?
\end{enumerate}
\end{tasksbox}
